\section{Conditional sampling from the nested tables}\label{sec:sampling}

The relations will be tested against an action chosen after all the tables
have been sampled.  We therefore prove statements about total variation of
finite distributions on table matches.  These statements do not refer to a
target metric space.  In particular, they will apply simultaneously to all
uniformly bounded functions on the relevant assignments, including functions obtained
from a subsequently chosen action.

\subsection{The master pattern and its critical exponents}

Recall the master pattern $\mathcal M$ from Section~\ref{sec:model},
and let $\mathcal F$ be its field set.  Its required entries are the
seven fan entries, the top-mark entry, and one entry from each of the
$90$ four-letter tables.  In each four-letter entry its first two fields are
identified with the corresponding core fields; its other two fields are
new and are used in no other entry.  The six-letter table is not part of
$\mathcal M$.  Thus $|\mathcal F|=212$.  Write $\mathcal R$ for the set of
its $98$ array entries and $E_\rho\subseteq\mathcal F$ for the field
support of $\rho\in\mathcal R$.  Each $\rho$ belongs to a different
Bernoulli array.  The order of its coordinates is fixed throughout.

A \emph{match} is a map $\sigma:\mathcal F\to S_n$ for which every
required array entry is present.  No injectivity is required of a single
match.  We write $\Omega_n$ for the random set of matches and, when it is
nonempty, write $\mu_N$ for its uniform probability measure.

Recall from Definition~\ref{def:admissible} that a set of fields is
admissible if it consists either of at most three core fields, at most
two in each fan, or of at most two fields of one square decoration.
The loose field belongs to no fan.
\begin{definition}\label{sample:def-test}
A \emph{conditional test} is a triple $(I,e,t)$ with
$I\subseteq\mathcal F$, $e\in\mathcal F\setminus I$,
$I\cup\{e\}$ admissible, and $t$ a positive integer.
\end{definition}

For $\theta=(\alpha,\beta,\delta)$ define $c_\rho(\theta)$ so that
the inclusion probability of entry $\rho$ is $N^{-c_\rho(\theta)}$.
Its value and its critical value $c_\rho=c_\rho(0)$ are
\begin{equation}\label{sample:eq-costs}
\begin{array}{c|c|c}
\text{entry}&c_\rho(\theta)&c_\rho\\ \hline
\text{fan }i&k_i-2+\alpha&k_i-2\\
\text{top mark}&12+\beta-7\alpha&12\\
\text{four-letter entry}&2+\delta&2.
\end{array}
\end{equation}
For a set of fields $W$ put
\begin{equation}\label{sample:eq-h}
h_\theta(W)=|W|-\sum_{\rho:E_\rho\subseteq W}c_\rho(\theta),
\qquad h(W)=h_0(W).
\end{equation}
Thus $N^{h_\theta(W)}$ is the expected number of assignments to $W$
for which all required entries supported in $W$ are present.
The function $h$ is submodular.  Indeed, for each fixed $E_\rho$,
the indicator of $E_\rho\subseteq W$ is supermodular, and all
$c_\rho$ are nonnegative.  Notice also that
\begin{equation}\label{sample:eq-master-mean}
h(\mathcal F)=212-(17+12+180)=3,
\qquad
\E|\Omega_n|=N^{h_\theta(\mathcal F)}
             =N^{3-\beta-90\delta}.
\end{equation}
The expectation is exact: each assignment uses one entry from each of
the distinct arrays, so all of its inclusion requirements are independent.

\begin{lemma}\label{sample:lem-budget}
For every $W\subseteq\mathcal F$,
\begin{equation}\label{sample:eq-budget}
h(W)\geq\min\{2,|W|\}.
\end{equation}
If $S$ is admissible and $S\subseteq W$, then
\begin{equation}\label{sample:eq-admissible-budget}
h(W)\geq |S|.
\end{equation}
\end{lemma}
\begin{proof}
First consider only the core part of $W$.  Let $a_i$ be its number of
fields in $F_i$, and let $\ell\in\{0,1\}$ indicate the loose field.
Before charging the top mark, the exponent is
\[
\sum_i\bigl(a_i-(k_i-2)\mathbf 1_{\{a_i=k_i\}}\bigr)+\ell
\ \geq\ \sum_i\min\{a_i,2\}+\ell.
\]
The top-mark cost is charged only for the complete core.  For that set
the exponent is $32-17-12=3$.  Consequently every core subset has
exponent at least the size of any admissible core subset it contains,
and at least the smaller of two and its own cardinality.

Now add the fresh fields of the four-letter entries.  Adding one fresh
field of an entry adds one to the exponent.  Adding both adds two and
can charge the cost two only if that entry's two core fields are also
present.  The net contribution of each pair of fresh fields is therefore
nonnegative.  This proves Equation~\eqref{sample:eq-admissible-budget}
when $S$ consists of core fields, and proves
Equation~\eqref{sample:eq-budget} when $W$ has at least two core fields.
If $W$ has fewer than two core fields, no four-letter entry is charged,
and no fan or top entry is charged either.  In that case $h(W)=|W|$.
This proves Equation~\eqref{sample:eq-budget} in all cases.  Finally,
an admissible subset of one four-letter entry has size at most two, so
its required bound follows from Equation~\eqref{sample:eq-budget}.
\end{proof}

Fix a conditional test $(I,e,t)$.  To couple $t$ matches agreeing on
$I$, we will require them to agree on a possibly larger set $U$.
The purpose of this enlargement is to make every further extension
increase the critical exponent.  For example, if $I$ consists of two
fields in one fan $F_i$, then completing that fan leaves the exponent
unchanged:
\[
h(I)=2=h(E(F_i)).
\]
Once the fan fields have been included, adjoining both fresh fields of
any square whose two core fields lie in that fan also leaves the exponent
unchanged.  These are examples of the enlargements that the choice of
$U$ must absorb.

In general, $|I|\leq2$, and every entry support has at least three
fields, so $h(I)=|I|$.  Choose an inclusion-maximal set
$U\supseteq I$ such that
\begin{equation}\label{sample:eq-closure}
h(U)=|I|.
\end{equation}
Such a set exists because the field set is finite.  The next two
properties are the reason for this choice:
\begin{equation}\label{sample:eq-closure-gap}
e\notin U,
\qquad
h(U\cup V)-h(U)\geq1
\quad\text{for every nonempty }V\subseteq\mathcal F\setminus U.
\end{equation}
The first follows from Equation~\eqref{sample:eq-admissible-budget}
applied to $I\cup\{e\}$.  For the second, every superset of $U$
has exponent at least $|I|$ by the same lemma applied to $I$.
Maximality rules out equality for a strict superset, and all critical
exponents are integers.

\subsection{Stars and their overlap estimates}

Form a star $\mathcal M_U^{(t)}$ from $t$ ordered copies of
$\mathcal M$ by identifying all copies of each field in $U$.
Entries supported entirely in $U$ are also identified.  Each other
entry retains one copy in each branch.  We require assignments to this
star to be \emph{injective in each field}: for each
$s\in\mathcal F\setminus U$, its $t$ branch values must be pairwise
distinct.  Values in different fields may coincide.  All fields continue
to take values in the same alphabet $S_n$.

This injectivity has two consequences.  First, copies of an entry not
supported in $U$ are distinct tuples in their array, since some one of
their coordinates lies outside $U$.  Second, when one branch is fixed,
an entry in another branch that is not supported in $U$ cannot equal
the corresponding entry of the fixed branch.  Thus all random variables
counted below are products of distinct Bernoulli trials, after explicitly
removing any trials on which we have conditioned.

We give the overlap count in detail because the coordinates are not
disjoint letter domains.  An \emph{overlap type} between two star
assignments specifies precisely which entries of each array coincide
across the two stars.  There are finitely many such types for fixed $t$.
Each equality of array entries equates corresponding coordinates.  For
a given field $s\notin U$, injectivity implies that these equalities
give a partial matching between the $t$ branch variables of the first
star and those of the second.  Different arrays may use different
partial matchings; their union must still be a partial matching whenever
the overlap type is realizable.  Otherwise two different branch values
of the same field would be forced equal.  For $s\in U$ there is just
one variable in each star.  In particular, an equality never compares
two different coordinate fields, even if their numerical letter values
happen to be equal.

Here is the resulting counting rule.  Reveal the first assignment.
Suppose the equalities of an overlap type prescribe the values of $q$
distinct, previously unpinned variables in the second assignment.
For each field with unpinned variables, there are a bounded number of
those variables, and they are chosen uniformly without replacement,
possibly avoiding one already pinned value.  A field whose variables
are all pinned contributes a factor one.  Prescribing $q_s$ of them has probability at most
$C_t N^{-q_s}$: the number of remaining choices is at most
$N^{r-q_s}$, whereas the number of unrestricted choices is a falling
factorial $(N-a)_r$, with $r\leq t$ and $a\in\{0,1\}$.
Multiplying over the finitely many fields gives
\begin{equation}\label{sample:eq-assignment-penalty}
\text{fraction of pairs having the specified overlap}
\ \leq\ C_tN^{-q}.
\end{equation}
Incompatible specifications simply contribute zero.

If the coincident, unconditioned entries have total actual cost $C$,
their joint inclusion probability divided by the product of the two
individual inclusion probabilities is $N^C$.  All other trials are
independent.  Consequently this overlap type contributes at most
\begin{equation}\label{sample:eq-variance-rule}
C_t N^{-q+C}
\end{equation}
to the relative variance of the count in question.  Pairs with no
coincident unconditioned entry have covariance zero, including when
they have coincident values that do not complete any common entry.
To justify the relative-variance normalization directly, all candidate
assignments in any one of our counting problems have the same
probability of inclusion, and Equation~\eqref{sample:eq-assignment-penalty}
bounds the fraction of pairs of such candidates.  Summing their
covariances and dividing by the square of the expectation gives
Equation~\eqref{sample:eq-variance-rule}, summed over the finitely many
nonempty overlap types.

We first apply this rule to a single copy of $\mathcal M$.  If $W$ is
the set of fields belonging to its shared entries, then $q=|W|$ and
the recovered cost is at most the sum of the costs of all entries
supported in $W$.  Its critical excess is therefore at least
$h(W)>0$.  For sufficiently small actual parameters, every one of
these finitely many positive excesses stays positive.  It follows that
\begin{equation}\label{sample:eq-single-concentration}
\frac{\operatorname{Var}|\Omega_n|}{(\E|\Omega_n|)^2}=o(1).
\end{equation}
In particular the master count is asymptotic to its expectation in
probability.  We always choose the parameters small enough that the
exponent in Equation~\eqref{sample:eq-master-mean} is positive.

For the two star counts, write
\[
C_U(\theta)=\sum_{\rho:E_\rho\subseteq U}c_\rho(\theta),
\qquad
C_{\mathrm{out}}(\theta)
 =\sum_{\rho:E_\rho\not\subseteq U}c_\rho(\theta),
\qquad d=|\mathcal F\setminus U|.
\]
Let $Z_N$ be the number of present injective stars.  With
$(N)_t=N(N-1)\cdots(N-t+1)$,
\begin{equation}\label{sample:eq-star-mean}
\E Z_N
 =N^{|U|}(N)_t^d
  N^{-C_U(\theta)-tC_{\mathrm{out}}(\theta)}.
\end{equation}

\begin{lemma}\label{sample:lem-conditional-counts}
For each fixed conditional test, sufficiently small actual parameters
have the following properties as $N\to\infty$.
\begin{enumerate}
\item Fix any assignment $\sigma:\mathcal F\to S_n$, place it in a
specified branch, and condition on it being a match.  Let $A_\sigma$
count its extensions to a present injective star.  Its conditional
mean is independent of $\sigma$ and equals
\begin{equation}\label{sample:eq-extension-mean}
a_N=((N-1)_{t-1})^d
       N^{-(t-1)C_{\mathrm{out}}(\theta)}.
\end{equation}
Uniformly in $\sigma$, its conditional relative variance is $o(1)$.
\item Fix any ordered list $b=(b_1,\ldots,b_t)$ of distinct letters
as the values of field $e$ in the branches.  Let $B_b$ count the present
injective stars having those values.  Its mean is independent of $b$
and equals
\begin{equation}\label{sample:eq-leaf-mean}
b_N=N^{|U|}(N)_t^{d-1}
       N^{-C_U(\theta)-tC_{\mathrm{out}}(\theta)}.
\end{equation}
Uniformly in $b$, its relative variance is $o(1)$.
\end{enumerate}
\end{lemma}
\begin{proof}
For the first assertion, condition only on the Bernoulli trials used by
$\sigma$.  All other trials retain their original independent laws.
Entries supported in $U$ have thereby been fixed present, as have the
other entries of the specified branch.  Every other branch uses new
entries distinct from those conditioned trials.  For each field outside
$U$, its other $t-1$ values are distinct and avoid the value in
$\sigma$, giving $(N-1)_{t-1}$ choices.  Each extra branch requires
all the entries not supported in $U$, giving exactly
Equation~\eqref{sample:eq-extension-mean}.

Consider an overlap of the unconditioned entries of two such extensions.
In extra branch $j$ of the first extension let $V_j$ consist of the
outside-$U$ fields involved in its shared entries.  The assignment
penalty in Equation~\eqref{sample:eq-assignment-penalty} is
$\sum_j|V_j|$: all $U$-values are fixed, whereas every involved
outside-$U$ variable is free and, by injectivity, cannot match a value
in the pinned branch.  Charge each shared entry to its branch in the
first extension.  The recovered critical cost there is at most
\[
\sum_{\rho:E_\rho\subseteq U\cup V_j,\ E_\rho\not\subseteq U}
c_\rho.
\]
Thus the critical excess of assignment penalty over recovered cost is
at least
\[
\sum_{j:V_j\ne\varnothing}
 \bigl[h(U\cup V_j)-h(U)\bigr]\geq1
\]
for every nonempty overlap, by
Equation~\eqref{sample:eq-closure-gap}.  The finitely many bounds remain
strict for small actual parameters.  Equations~\eqref{sample:eq-assignment-penalty}
and~\eqref{sample:eq-variance-rule} prove the conditional relative
variance claim, uniformly in the values of $\sigma$.  If $t=1$, the
extension count is identically one and there is no variance to estimate.

For the second assertion, no Bernoulli trials are conditioned on;
only letter values are pinned.  Since $e\notin U$, its $t$ pinned
values satisfy precisely its field-injectivity constraint.  The other
fields have $N^{|U|}(N)_t^{d-1}$ choices, proving
Equation~\eqref{sample:eq-leaf-mean}.

For an overlap let $W$ be the set of variables in the first star
belonging to shared entries.  Write $A=W\cap U$, and let $V_j$ be
its outside-$U$ fields in branch $j$.  Let $m(W)$ be the number of its
$e$-variables.  In a realizable overlap an $e$-variable can only match
the variable with the same pinned value; those equalities impose no
further assignment penalty.  All other involved variables do impose
one.  Thus $q=|W|-m(W)$ in
Equation~\eqref{sample:eq-assignment-penalty}.

Define $h_\star(W)$ using the critical entry costs in the star, counting
entries entirely in $U$ once and the other entries separately in each
branch.  The recovered cost is at most the cost of all star entries
supported in $W$.  Therefore the critical excess is at least
\begin{equation}\label{sample:eq-leaf-excess}
h_\star(W)-m(W).
\end{equation}
If $A\ne\varnothing$, separating entries entirely in $U$ from the
entries in each branch gives
\begin{align}
h_\star(W)
&\geq h(A)+
 \sum_{j:V_j\ne\varnothing}
       \bigl[h(U\cup V_j)-h(U)\bigr].
\label{sample:eq-star-submodular}
\end{align}
Indeed the exact contribution of branch $j$ is $|V_j|$ minus the
costs of entries supported in $A\cup V_j$ but not entirely in $U$;
replacing $A$ by $U$ can only increase those costs.  This proves the
displayed inequality directly.  Its first term is at least one by
Lemma~\ref{sample:lem-budget}, and each nonempty summand is at least
one by Equation~\eqref{sample:eq-closure-gap}.  Each such branch
contains at most one pinned $e$-variable.  Hence the excess in
Equation~\eqref{sample:eq-leaf-excess} is at least one.

If $A=\varnothing$, no shared entry meets $U$, and entries in different
branches use disjoint variables.  Consequently
$h_\star(W)=\sum_j h(V_j)$.  An active branch contains the full
support of at least one shared entry.  Every entry support has at least
three fields, so $|V_j|\geq3$ in such a branch and
Lemma~\ref{sample:lem-budget} gives $h(V_j)\geq2$.
Subtracting at most one pinned leaf per branch again leaves a positive
excess.  This argument is applied to supports of shared entries, not
to arbitrary nonempty field subsets.

All nonempty overlap types thus have critical excess at least one.
Perturbing to sufficiently small actual parameters preserves a positive
excess for every type.  Equation~\eqref{sample:eq-variance-rule}, summed
over those types, proves the second relative variance assertion.
The bounds are uniform in the pinned letters because all assignment
counts and Bernoulli probabilities used above are independent of their
particular values.
\end{proof}

\subsection{Normalization to probability couplings}

For probability measures $\lambda,\lambda'$ on a finite set we use
\[
d_{\mathrm{TV}}(\lambda,\lambda')
 =\frac12\sum_z|\lambda(z)-\lambda'(z)|.
\]
All convergence statements about the random tables below are in
probability as $N\to\infty$; equivalently one can choose deterministic
error bounds tending to zero for which the corresponding assertions
hold with probability tending to one.

\begin{lemma}\label{sample:lem-coupling}
For each fixed conditional test and sufficiently small actual
parameters, with probability tending to one there is a probability
law on $t$ matches $(\sigma_1,\ldots,\sigma_t)$ such that:
\begin{enumerate}
\item their restrictions to $I$ agree identically;
\item every branch marginal has total-variation distance $o(1)$ from
$\mu_N$;
\item the law of $(\sigma_1(e),\ldots,\sigma_t(e))$ has
total-variation distance $o(1)$ from the uniform product law on $S_n^t$.
\end{enumerate}
\end{lemma}
\begin{proof}
Use the uniform law on present injective stars.  We verify that its
normalization is nonzero and has the stated properties.

Put $L_N=|\Omega_n|$ and $\lambda_N=\E L_N$.
For a specified branch let
\[
D_N=\sum_{\sigma\in\Omega_n}|A_\sigma-a_N|.
\]
The conditional variance bound of
Lemma~\ref{sample:lem-conditional-counts}, followed by the
Cauchy--Schwarz inequality, gives uniformly in each possible assignment
$\sigma$
\[
\E\bigl[|A_\sigma-a_N|\mid \sigma\in\Omega_n\bigr]=o(a_N).
\]
Multiplying by the probability that $\sigma$ is present and summing
over all assignments shows
$\E D_N=o(\lambda_Na_N)$.  Markov's inequality yields
$D_N=o(\lambda_Na_N)$ in probability.  By
Equation~\eqref{sample:eq-single-concentration},
$L_N/\lambda_N\to1$ in probability.  Since
\[
Z_N=\sum_{\sigma\in\Omega_n}A_\sigma,
\qquad |Z_N-L_Na_N|\leq D_N,
\]
we have $Z_N/(L_Na_N)\to1$ in probability, and in particular
$Z_N>0$ with probability tending to one.  Moreover,
\begin{align*}
\sum_{\sigma\in\Omega_n}
 \left|\frac{A_\sigma}{Z_N}-\frac1{L_N}\right|
&\leq \frac{D_N}{Z_N}
      +\left|\frac{L_Na_N}{Z_N}-1\right|
\longrightarrow0.
\end{align*}
The same argument applies to each of the finitely many branches.
This proves the marginal assertion.  Agreement on $I$ is exact because
the branches agree on $U\supseteq I$.

For the leaves, let $\mathcal B_N$ be the set of ordered distinct
letter lists of length $t$, so $|\mathcal B_N|=(N)_t$.  By the second
part of Lemma~\ref{sample:lem-conditional-counts},
\[
\E\sum_{b\in\mathcal B_N}|B_b-b_N|
       =o((N)_t b_N).
\]
Markov's inequality and $\sum_bB_b=Z_N$ imply that the total variation
between the leaf distribution $B_b/Z_N$ and the uniform distribution
on $\mathcal B_N$ tends to zero, by the same displayed normalization
estimate with $(N)_t,b_N$ in place of $L_N,a_N$.
Finally, an iid uniform list in $S_n^t$ has a repeated value with
probability at most $\binom t2/N$.  Its distribution conditioned on
distinctness is uniform on $\mathcal B_N$.  This proves the final
assertion.
\end{proof}

\subsection{Finite tests, parameter choices, and the six-letter table}

The two marginal conclusions of the coupling serve different purposes.
The branch comparison transfers averages of bounded tests between the
coupled branches and a uniform match.  In
Lemma~\ref{lem:conditional-test}, a pointwise norm estimate bounds the
average over the $t$ branches by a function of the $e$-labels alone; the common
$I$-labels enter only through a uniform bound.  The comparison with
independent uniform letters then controls that function of the leaves.
Thus the argument needs no independence between the common $I$-tuple
and the test labels.

All critical overlap margins above are at least one.  Each shared trial
is charged once to its occurrence in the first star, which has at most
$t$ copies of each array entry.  Since $\beta>7\alpha$, all cost
increments are positive.  Replacing critical costs by actual costs
therefore increases any recovered-cost sum by at most
\[
t\bigl[7\alpha+(\beta-7\alpha)+90\delta\bigr]
=t(\beta+90\delta).
\]
Conditioning removes trials and can only decrease this bound; a single
master match is covered by $t=1$.  For a finite family of tests, let
$t_{\max}$ be the largest requested value of $t$, or one if the family
is empty.  Choosing $t_{\max}(\beta+90\delta)<1/2$ preserves every
overlap margin at a value at least $1/2$, including the singleton
tests used below.  This is compatible with $\beta>7\alpha>0$: take
$\beta=8\alpha$ and then take $\alpha,\delta$ sufficiently small.
The parameters are held fixed while $N$ tends to infinity.  In
particular, none of the variance arguments asserts uniformity for
unbounded $t$ at a single fixed choice of parameters.

Let $\mathcal D_N\subseteq S_n^6$ be the separate six-letter table.
For positions $1\leq r<q\leq6$ and letters $u,v$, write
\[
Q_{rq}(u,v)
 =\#\{(s_1,\ldots,s_6)\in\mathcal D_N:s_r=u, s_q=v\}.
\]
Each such count is a sum of $N^4$ independent Bernoulli variables
of parameter $N^{-3-\delta}$, and so has mean $N^{1-\delta}$.
For fixed $0<\delta<1$, put
$\varepsilon_N=N^{-(1-\delta)/4}$.  The binomial Chernoff bound gives
\[
\Prob\bigl(|Q_{rq}(u,v)-N^{1-\delta}|>
       \varepsilon_NN^{1-\delta}\bigr)
 \leq 2\exp\bigl(-\varepsilon_N^2N^{1-\delta}/3\bigr).
\]
A union bound over the $15N^2$ choices of positions and letters
therefore proves
\begin{equation}\label{sample:eq-six-pair}
\max_{r<q}\max_{u,v\in S_n}
 \left|\frac{Q_{rq}(u,v)}{N^{1-\delta}}-1\right|=o(1)
\end{equation}
with probability tending to one.  No independence between these
different counts is needed.

We collect the conclusions in the form used below.
\begin{proposition}\label{prop:sampling}
For any finite family $\mathcal T$ of conditional tests, there is a
neighborhood of the critical parameters such that every sufficiently
small fixed positive $(\alpha,\beta,\delta)$ in that neighborhood
with $\beta>7\alpha$ has the following properties as $N\to\infty$.
With probability tending to one, the master pattern has a nonempty
set of matches, its count is
\[
|\Omega_n|=(1+o(1))N^{3-\beta-90\delta},
\]
every single-field marginal of its uniform law $\mu_N$ is $o(1)$
in total variation from the uniform law on $S_n$, and every test in
$\mathcal T$ has the coupling of Lemma~\ref{sample:lem-coupling}.
The six-letter counts also satisfy Equation~\eqref{sample:eq-six-pair}.
All these properties concern only the tables and are independent of
any choice of action or metric space.
\end{proposition}
\begin{proof}
Adjoin to $\mathcal T$ the finitely many tests
$(\varnothing,e,1)$, one for each field $e$.  For such a test
Equation~\eqref{sample:eq-budget} forces $U=\varnothing$, and its
one-branch star is exactly a master match.  The leaf assertion in
Lemma~\ref{sample:lem-coupling} therefore gives single-field
uniformity.  Choose parameters satisfying all the finite overlap
margins just described, $\beta+90\delta<3$, and $\delta<1$.
Equation~\eqref{sample:eq-single-concentration},
Lemma~\ref{sample:lem-coupling}, and
Equation~\eqref{sample:eq-six-pair}, combined over this finite set of
events, prove the proposition.
\end{proof}

\begin{remark}\label{sample:rem-diagonal}
Proposition~\ref{prop:sampling} is compatible with
Theorem~\ref{thm:diagram}: for each finite family of tests, choose
positive parameters sufficiently small for that family and the theorem,
then the diagram fragment threshold, and finally $N$.
The diagram and sampling events then hold simultaneously with probability
tending to one.  Lemma~\ref{lem:completion-diagonal} constructs the
deterministic diagonal sequence used below, with parameters tending to
zero and every fixed test valid in the limit.
\end{remark}
